\documentclass[10pt]{article}

\usepackage[utf8]{inputenc}

\usepackage[T1]{fontenc}

\usepackage{amsmath}

\usepackage{amsfonts}

\usepackage{amssymb}

\usepackage[version=4]{mhchem}

\usepackage{stmaryrd}

\usepackage{mathrsfs}



\begin{document}

Лит.: [1] Бе рез и н Ф.А., Ш уб и н М.А., Уравнение Шредингера, М., 1983. М. А. Шубин.



KВАНТОВАНИЕ IPОСТРАНСТВА-ВРЕМЕНИ - общее название обобщений квантовой теории поля, основанных на гипотезе о существовании фундаментальной длины как одной из универсальных физических постоянных (наряду с $\hbar$ и c ). Ближайшая цель таких обобщений - освобождение от расходимостей, появляющихся в традиционной квантовой теории поля. (СМ. таКже Нелокальная квантовая теория поля.)



При построении теории, описывающей явления микромира, классич. представления о пространстве и времени, в частности представление о принципиальной возможности сколь угодно точного измерения расстояний (длин) и промежутков времени, были без каких-либо изменений перенесены в новую область. Введение фундаментальной (минимальной) длины $l$ соответствует предположению, что измерение малых расстояний возможно лишь с ограниченной точностью порядка $l$ (и времени - с точностью порядка $l / c$ ).



Существует несколько способов введения фундаментальной длины. Один из них связан с переходом от непрерывных значений координат к дискретным величинам (наподобие правил квантования Бора в первоначальной теории атома), другие - с заменой координат и времени на некоммутирующие между собой операторы (наподобие операторов координаты $\hat{x}$ и импульса $\hat{p}$ в квантовой механике), вследствие чего координаты не могут иметь точных значений в данный момент времени. Вид операторов подбирается так, чтобы средние значения координат могли принимать лишь значения, кратные фундаментальной длине $l$. Во всех вариантах введение минимальной длины исключает существование волн с длиной $\lambda<l$, то есть как раз тех квантов бесконечно большой энергии $\mathscr{E}=2 \pi \hbar c / \lambda$, с к-рыми связано появление ультрафиолетовых расходимостей. Однако введение фундаментальной длины, по-видимому, не устраняет основного противоречия квантовой теорий поля, связанного с возможностью неограниченного роста эффективного заряда с уменьшением расстояния. Все же такой пересмотр может оказаться необходимым.



Лum.: [1] Марко в М.А., Гипероны и К-мезоны, М., 1958; [2] Блохи нце в Д. И., Пространство и время в микромире, М., 1970.



КВАНТОВАНИЯ УСЛОвИй ПОПРАВКИ - числовые топологические характеристики лагранжевых подмногообразий (одномерные классы когомологий), возникаюшие при построении квазиклассической асимптотики спектральных серий и добавляюшиеся, помимо индекса Маслова, в Бора-Зоммерфельда условия квантования за счет деформации симплектической структуры или при замене стандартного фазового пространства нетривиальным расслоением (со связностью). Такие поправки появляются в уравнениях с операторными коэффициентами и учитываются по теории (см. [1]), напр., в условии квантования с периодич. магнитным полем (см. [2], [3]). Аналогичную структуру имеют поправки к условиям квантования изотропных подмногообразий с комплексным ростком. К. у. П. дают также деформации скобок Пуассона между симметриями или динамич. переменными (см. [4]).



Лит.: [1] М а с л о в В. П., Теория возмущений и асимптотические методы, M., 1965; [2] B e r ry M. V., «Proc. Roy. Soc. London», 1984, v. A392, p. 45-57; [3] Бу сл а е в В. С., «Успехи математич. наук», 1987, т. 42, в. 6, с. 77-98; [4] В о р б бе в Ю. М., Ка расе в М. В., «Доклады АН СССР», 1987, т. 297, № 6, с. 1294-98. Ю. М. Воробьев. КВАНТОвАНИЯ УслОвМЯ Маслова - см. Маслова канонический оператор.



KBAHTOBAHHOE MHOTOOEPAЗИE - см. Маслова канонический оператор.



KВАНТОВАЯ ГЕОМЕТРОДИНАМИКА - см. Квантовая теория гравитации.



KBAНTOBAЯ ГPУППА - по сушеству, то же, что Хопфа алгебра. Группе $G$ соответствует алгебра Хопфа $A$, состоя-



244 KBAHTOBAHИE щая из функций на $G$, где умножение поточечное, а коумножение $\Delta: A \rightarrow A \otimes A$ переводит $f \in A$ в функцию $\varphi$ на $G \times G$ :



\[

\varphi(g, h)=f(g h)

\]



(функции на $G \times G$ принацлежат $A \otimes A$ ). Так получается биекция между обычными группами и коммутативными алгебрами Хопфа. Произвольную (не обязательно коммутативную) алгебру Хопфа $A$ естественно рассматривать как алгебру функций на К. г. Напр., алгебру $A$ с образуюшими $t_{i j}, 1 \leqslant i, j \leqslant 2$, соотношениями



\[

\begin{gathered}

{\left[t_{12}, t_{21}\right]=0,} \\

{\left[t_{11}, t_{22}\right]=\left(e^{h}-e^{-h}\right) t_{12} t_{21}, t_{i 1}, t_{i 2}=e^{h} t_{i 2} t_{i 1},} \\

t_{1 j} t_{2 j}=e^{h} t_{2 j} t_{1 j}, t_{11} t_{22}-e^{h} t_{12} t_{21}=1,

\end{gathered}

\]



где $h$ - постоянная Планка, и коумножением



\[

\Delta\left(t_{i j}\right)=t_{i 1} \otimes t_{1 j}+t_{i 2} \otimes t_{2 j}

\]



можно рассматривать как алгебру функций на квантованной группе $S L(2)$. К.г. (в частности, квантовые аналоги группы функций одного переменного со значениями в простой группе Ли) естественным образом возникают в квантовом методе обратной задачи.



См. также Кольцевая группа.



Лит.: [1] Д р и нфе льд В.Г., «Зап. науч. семинаров ЛОМИ», 1986, т. 155, с. 18-49; [2] Решетихин Н. Ю., Тахтаджян Л. А., Фа дде е в Л. Д., “Алгебра и анализ", 1989, т. 1, в. 1, с. 178-206.



В. Г. Дринфельд.



KВАНТОВАЯ ДИНАМИЧЕСКАЯ ПОЛУГРУППА - см. Динамическая полугруппа.



KВАНТОВАЯ ДИНАМИЧЕСКАЯ СИСТЕМА в Уз ком с м ы с е - это $C^{*}$-алгебра вместе с двухпараметрическим семейством $\left\{\Theta_{t}^{S}\right\}$ согласованных относительно композиции $\Theta_{r}^{t} \Theta_{t}^{S}=\Theta_{r}^{s}, \quad r<t<s, \quad$ эндоморфизмов $\Theta_{t}^{s}, \quad t<s, \quad$ имеющих обычно пространственный вид $\Theta_{t}^{S}(a)=V_{t}^{s} \stackrel{*}{t} V_{t}^{s}$, где $V_{t}^{s}: H \rightarrow H, t<s,-$ изометрические (унитарные) пропагаторы, удовлетворяюшие условию



\[

V_{t}^{S} V_{r}^{t}=V_{r}^{S}, r<t<s

\]



Квантовая динамическая система в широком с мы с л е описывается (см. [1], [2]) согласованным семейством $\left\{\Theta_{t}^{s} \mid t \leqslant s\right\}$ линейных вполне положительных отображений $\Theta_{t}^{s}: \mathfrak{A} \rightarrow \mathfrak{A}$, удовлетворяющих условию нормировки $\Theta_{t}^{S}(1)=1$ относительно единицы $1 \in \mathfrak{A}$. Отображения $\left\{\Theta_{t}^{S}\right\}$ определяют квантовую регрессию хронологически-упорядоченных корреляционных форм, с к-рыми теорема реконструкции (см. [1]) связывает однозначно минимальный квантовый марковский процесс $\pi_{t}: \mathfrak{A} \rightarrow \mathscr{B}(H)$, удовлетворяющий условию динамичности



\[

\pi_{t} \circ \Theta_{t}^{s}(a)=E_{t} \pi_{s}(a) E_{t}, a \in \mathfrak{A},

\]



относительно ортопроекторов $E_{t}=\pi_{t}(1)$.



В стационарном случае $\Theta_{t}^{S}=\Theta(s-t)$ динамич. отображения $\Theta(s)=\Theta_{0}^{S}$ образуют полугруппу $\Theta(s) \Theta(r)=\Theta(s+r)$, структура генераторов к-рой установлена в равномерном случае (cM. [2]).



Лuт.: [1] Белавкин В. П., «Теоретич. и математич. физика», 1985 , т. 62, № 3, c. 409-31; [2] Lind blad G., «Comm. Math. Phys.», 1976, v. 48, № 1, p. 119-30. В. П. Белавкин. КВАНТОВАЯ КОСМОЛОГИЯ - см. Квантовая теория гравитации.



KВАНТOBAЯ ЛОГИКА - множество высказываний о всевозможных событиях в физической системе, составляющих ее микроскопическое описание (см. [1]). В отличие от классич. логики, подчиняющейся законам булевой алгебры, высказывания К. л. образуют (см. [2]) более общую ортомодулярную структуру (решетку) относительно логич. порядка « $a$ влечет $b »$, обозначаемого $a \leqslant b$, и инволютивной операции отрицания $a \mapsto \bar{a}, \overline{\bar{a}}=a$, обращающей этот порядок:





\end{document}